\documentclass[10pt]{article}

\def\Type{LessonCaelan}
\def\TypeNum{ End of Semester Reflection} %Lesson #
\def\TypeTitle{ } %Lesson Title
\def\Semester{Fall 2021} %Not Used 
\def\Course{MATH 1190}%Course number and name
\def\Targets{   \target{P1}, \target{P2} }

\usepackage{preambleDJ} 

%%%%%%%%%%%% BEGIN DOCUMENT %%%%%%%%%%%%%%%
%%%%%%%%%%%% BEGIN DOCUMENT %%%%%%%%%%%%%%%
%%%%%%%%%%%% BEGIN DOCUMENT %%%%%%%%%%%%%%%
%%%%%%%%%%%% BEGIN DOCUMENT %%%%%%%%%%%%%%%
%%%%%%%%%%%% BEGIN DOCUMENT %%%%%%%%%%%%%%%

\begin{document}
\pagestyle{otherpages}
\thispagestyle{firstpage}


\vs{0}


%



\section*{Reflection on course material}

\newcommand{\contentquestion}[1]{
    \begin{enumerate}
        \item On a scale of 1 to 5, rate your understanding of #1:\\[.6cm]
            before taking this course:\rule[-8pt]{1in}{0.4pt}\hspace{2cm}
            at this point in the course:\rule[-8pt]{1in}{0.4pt}
            
        \item Were there any specific moments where your understanding of #1 ``clicked", or study resources that helped you understand #1 better?
        \vspace*{3cm}

        \item Is there anything about #1 that is still confusing to you?
        \vspace*{3cm}

        \item Did you struggle to keep up with the class during the #1 part of the course? If so, how did you respond?
        \vspace*{3cm}
    \end{enumerate}
}

In this activity, you will reflect on how you have grown in your understanding of the various math concepts covered in this course. When rating your understanding on a scale of 1 to 5, use 1 as the lowest (no understanding whatsoever) and 5 as the highest (able to solve novel problems and confidently explain this concept to others).

\begin{enumerate}
    \item \textbf{Limits:} The meaning of limits of functions, continuity, and evaluating limits using forms.
    \contentquestion{limits}
\newpage
    \item \textbf{Derivatives:} The meaning of the derivative, and computing derivatives of functions using a variety of rules.
    \contentquestion{derivatives}

\newpage
    \item \textbf{Optimization:} Critical numbers, concavity, and finding local and global extrema of functions.
    \contentquestion{optimization}
\newpage
    \item \textbf{Integrals:} The meaning of definite integrals, and computations with the Fundamental Theorem of Calculus.
    \contentquestion{integrals}
\newpage
    \item \textbf{Probability:} Discrete and continuous random variables, probability mass functions and probability density functions, and performing computations using these.
    \contentquestion{probability}
\end{enumerate}

\end{document}